\documentclass{article}
\usepackage[utf8]{inputenc}
\usepackage{amsmath}
\usepackage{amssymb}
\usepackage{graphicx}
\usepackage{hyperref}
\usepackage{tikz}
\usetikzlibrary{calc} % Needed for coordinate calculations


\title{Problem Solving Strategies - Chapter 3}
\author{Thomas Kidane}
\date{\today}

\begin{document}

\maketitle

\tableofcontents

\section{Introduction}
This might actually be one of the most important sections. I didn't fully understand the examples but hope that the 
problems will atleast teach me something.

\section{Problems}

\subsection{Problem 1}

Prove that there are at least $(2n-2)/3$ triangles among the $p_{n}$ parts of the plane in
Example \#1.

\textbf{Solution}

The first example is 

(a) Into how many parts at most is a plane cut by n lines? (b) Into how many
parts is space divided by n planes in general position?

$p_{n}$ is the max number of lines. $p_{n}=\binom{n}{0}+\binom{n}{1}+\binom{n}{2}$.

How can I uniquely determine a triangle? Could I determine it by the highest point at the top? Probably,
I could always reorient the plane enough so that there is only one point. 

What is needed to determine this max point? I would atleast 2 lines?

Wait, wait? What if I try to add interpretation to the equation? It seems like it is telling me that there are 
about $2n-2$ edge point(vetices, or corners) and that only a 3 make up the triangle?

Well are there really $2n-2$ intersection points? There should be about $\binom{n}{2}$ intersection points. Which 
comes out to be $n*(n-1)/2$ points. I guess this is 


Wait, this specifically specifies trianlges? I counted just inerseciton points, it could the case that there are 
lots of quadrilaterals. But how do I guarantee that I get the triangels?

Which set of lines could make a triangle? I neeed an extremal property? 

Actually let's make this interpretation, on each the lines there are 2 triangles, except one special line. 
Each traingle is counted 3 times, so we must divide by 3. 

Now how do I show that 2 triangles rest on each line, pick a side, then choose the intersection 
point that is closest to the lines. This must make a triangle, because there is no intersection point that is closer.

So now I need to prove that there are max 2 lines that don't have an intersection on one side. It should be the case that 
if there are more than 2 then I can turn that into a triangle. 

Lowk, am I cooked? I think I need to use an argument that it would violate the number of areas or somethigng but I don't know. 

And the book doesn't say anything else. Guess I am cooked. 

\subsection{Problem 2}

In the plane, $n$ lines are given ($n\geq 3$), no two of them parallel. Through every
intersection of two lines there passes at least an additional line. Prove that all lines
pass through one point.

\textbf{Solution}

Assume that there it isn't the case that there is only one point of intersection. Then take the furthest
intersection point to the left and then to the bottom. 


\subsection{Problem 5}

Does there exist a tetrahedron, so that every edge is the side of an obtuse angle of a face? 

\textbf{Solution}

I couldn't solve this one, but after taking a look at the answer it was so simple(just like all the problems).

So basically you take the longest edge in the tetrahedron, and you look at the look triangles it is a part of, there are 2.
Since it is the longest, the side opposite the edge must be the greatest angle, therefore it can't be the case that this edge is an edge for
an obtuse angle.

\subsection{Problem 6}

Prove that every convex polyhedron has at least two faces with the same number of sides. 

\textbf{Solution}

This is just an application of pigeohole principle, where you take the face with the maximal number of edges, call this number n. 
All the other faces must have between $[3,n]$. Since there atleast $n$ other faces there must atleast be one 1 repetition.

\subsection{Problem 7}

$(2n + 1)$ persons are placed in the plane so that their mutual distances are different. 
Then everybody shoots his nearest neighbor. 

Prove that

(a) at least one person survives;\\
(b) nobody is hit by more then five bullets;\\ (c) the paths of the bullets do not cross;\\
d) the set of segments formed by the bullet paths does not contain a closed polygon.\\

\textbf{Solution}

a) Take the two people, that are closest to each other, this people are point their guns at each other.

Then take the set of people that are not this, they can be grouped on whether they are pointing at one of this people. 

Make $S$ where S contains this two people that are point to each other, and the other people that are point to the 2 people. 
Define $f : S \mapsto S^{*}$, where $S^{*}=S \cup $\{ people point their guns at people in S\}. Continually apply $f$ to itself until the result stops growing.
The last person added to this set will survive. Since now one is point their guns at him/her.

b) Take person who is hit by the most number of bullets. Then that means he is the person that is 
closest to the most number of people at the same time.

Take the person that is closest to him. Wait lowk, I think this is six bullets, since I can inscribe a perfect hexagon in a circle and everbody will have equal distance to each other.
Actually after reading the mutual distance condition, then it is right. Also the book has a good angle argument, which I should have refined.

c) Yes, they don't cross because of the traingle inequality.

d) It can't because it can't contain a cycle. Should have explicitly shown the inequality.


My proof have no substance behind them. Fuck. Wait let there exist a closed polygon. Take the two closest people in this polygog, they are point their guns at each other. 
So now you have a chain, but it can't go back on itself. 

\subsection{Problem 8}

Rooks are placed on the $n \times n$ chessboard satisfying the following condition:

If the square (i, j) is free, then at least n rooks are on the $i^{th}$ row and $j^{th}$ column together.

Show that there are at least $n^2$/2 rooks on the board. 

\textbf{Solution}

Take the row/column that has the least number of rooks. Let this number be $x$, then there are atleast $max(n*x,n*(n-x))$, and the 
minimum of this expression is when $x=n/2$, therefore the minimum is $n^2/2$.

\subsection{Problem 9}

All plane sections of a solid are circles. Prove that the solid is a ball.


\textbf{Solution}

I don't know.

The shortest proof runs as follows. Consider the largest chord of the solid. Any
section through this chord is a circle whose diameter is the chord. Otherwise the
circle and the solid would have a larger chord. Thus the solid is a ball and one of its
diameters is the selected chord.
This proof is not complete. We did not prove that a longest chord exists. In fact, if
the surface of the solid did not belong to the solid, a longest chord would not exist.
So we assume that the solid is a closed and bounded set. Then we can apply the
theorem of Weierstrab: A continuous function defined on a closed and bounded set
always assumes its global maximum and minimum.
This theorem belongs to higher mathematics, but at the IMO you can use it. The proof
is not considered to have a gap if you cite the theorem. There are also elementary
proofs which are slightly longer (see HMO 1954).

\subsection{Problem 10}

A closed and bounded figure $\phi$ with the following property is given in a plane: Any
two points of $\phi$ can be connected by a half circle lying completely in $\phi$ . Find the
figure $\phi$ (West German proposal for IMO 1977).

\textbf{Solution}

I don't know but after looking at the solution from the book, I realize it is maybe to 
geometry flavour for me, so I have listed the full solution for the sake of completeness below.

Too lazy ;(

\begin{tikzpicture}
  % Define points
  \coordinate (A) at (0,0);
  \coordinate (B) at (5,0);

  % Draw the longest segment AB
  \draw[thick] (A) -- (B) node[midway, below] {$AB$};

  % Draw the semicircle s with diameter AB
  % Calculate the center of AB
  \coordinate (MidAB) at ($ (A)!0.5!(B) $);
  % Calculate the radius of the semicircle on AB
  \pgfmathsetmacro{\radiusAB}{0.5 * veclen(5,0)} % Distance from A to B is 5, radius is 2.5
  \draw (A) arc (180:0:\radiusAB);

  % Label the semicircle s
  \node at (MidAB) [above=0.5cm] {$s$}; % Position relative to MidAB

  % --- Option 1: Point C INSIDE semicircle s ---
  \coordinate (C) at (2,1); % C chosen to be inside the semicircle s
  \draw[fill=red] (C) circle (0.05) node[above right] {$C$};

  % Draw AC and BC for point C
  \draw[dashed] (A) -- (C);
  \draw[dashed] (B) -- (C);

  % Draw semicircle with diameter AC
  \coordinate (MidAC) at ($ (A)!0.5!(C) $);
  \pgfmathsetmacro{\radiusAC}{0.5 * veclen(C.x - A.x, C.y - A.y)}
  \draw[blue, dashed] (MidAC) arc (180:0:\radiusAC); % Arc from left to right

  % Draw semicircle with diameter BC
  \coordinate (MidBC) at ($ (B)!0.5!(C) $);
  \pgfmathsetmacro{\radiusBC}{0.5 * veclen(B.x - C.x, B.y - C.y)}
  \draw[green, dashed] (MidBC) arc (0:180:\radiusBC); % Arc from right to left


  % --- Option 2: Point C OUTSIDE semicircle s (uncomment to see this) ---
  % If you want to show C outside, comment out Option 1 and uncomment this:
  % \coordinate (C_out) at (6,2); % C chosen to be outside the semicircle s
  % \draw[fill=red] (C_out) circle (0.05) node[above right] {$C$};
  %
  % % Draw AC and BC for point C_out
  % \draw[dashed] (A) -- (C_out);
  % \draw[dashed] (B) -- (C_out);
  %
  % % Draw semicircle with diameter AC_out
  % \coordinate (MidAC_out) at ($ (A)!0.5!(C_out) $);
  % \pgfmathsetmacro{\radiusAC_out}{0.5 * veclen(C_out.x - A.x, C_out.y - A.y)}
  % \draw[blue, dashed] (MidAC_out) arc (180:0:\radiusAC_out);
  %
  % % Draw semicircle with diameter BC_out
  % \coordinate (MidBC_out) at ($ (B)!0.5!(C_out) $);
  % \pgfmathsetmacro{\radiusBC_out}{0.5 * veclen(B.x - C_out.x, B.y - C_out.y)}
  % \draw[green, dashed] (MidBC_out) arc (0:180:\radiusBC_out);

\end{tikzpicture}

\subsection{Problem 11}
Of $n$ points in space, no four lie in a plane. Some of the points are connected by
lines. We get a graph $G$ with $k$ edges.
\begin{enumerate}
    \item[(a)] If $G$ does not contain a triangle, then $k \leq \lfloor n^2/4 \rfloor$.
    \item[(b)] If $G$ does not contain a tetrahedron, then $k \leq \lfloor n^2/3 \rfloor$.
\end{enumerate}


\textbf{Solution}

Where do I even start? What does Triangle have to with edges? What does the 4 points not in a plane thing have to do with anything?
Well if there are 4 points in a plane, then There are 6 edges between them? What if they not in a plane then?
Actually you can always make a triangle with 3 or more points, so the condition is basically saying that the points are not coplanar.

What do I do with this information? How does this affect anyting though? What is the natural upper bound on the number of distinct edges?
Well that would be $\binom{n}{2}$. 

Holly shit, I reread the quesiton, I choose some subset of edges from the $\binom{n}{2}$, edges, and I call that Graph $G$. 
Now this graph doesn't have a triangel? So that are not cycles of length 3. That is what is it saying?

So what is the maximum of $E(G)$ for a graph $G$ with n nodes such that it doesn't have a cycle of length 3. 




\end{document}
